%% ---------------------------------------------------------------------------
\section{Finite element data generation with PyNite}\label{sec:data}

\subsection{Parametric frame family}
The study considers planar, multi-storey, multi-bay moment-resisting frames with rigid beam--column joints, a family that is routinely re-analysed during preliminary design and assessment. A frame is defined by the number of storeys $n_s$ and bays $n_b$, the bay widths $L_b$, the storey heights $h_s$, Young's modulus $E$, the second moments of area of the columns ($I_{c,s}$) and beams ($I_{b,s}$) in each storey, the corresponding cross-sectional areas $A = I/r^2$ obtained from a sampled radius of gyration $r$, the base condition (fixed or pinned), the uniformly distributed gravity load $w_{s,b}$ on each beam and a lateral force $F_s$ at each floor. All columns of a storey share one section and column stiffness may decrease with height through a taper factor. Units are kN and m throughout. The sampling ranges are listed in Table~\ref{tab:ranges}; $E$ spans concrete-like to steel-like values so that the surrogate cannot rely on a fixed material.

Three scenarios are generated (Fig.~\ref{fig:frames}):
\begin{itemize}
\item \emph{Training design space (ID).} $1\le n_s\le 5$, $1\le n_b\le 4$; the gravity load is uniform along each floor but varies between floors; the lateral profile is either uniform or inverted-triangular, in either direction.
\item \emph{Unseen geometries (OOD-G).} Frames with $6$--$8$ storeys and/or $5$--$6$ bays, and in half of the cases a tall first storey ($4.5$--$6.0$\,m), all outside the training ranges.
\item \emph{Unseen loading (OOD-L).} Training geometries subjected to pattern gravity loading, in which the live part of the load is randomly removed from individual bays, and to lateral profiles that were never seen in training (parabolic, or a single force at the roof).
\item \emph{Design-driven frames (DES).} 300 steel moment frames whose parameters are correlated as in practice instead of being sampled independently: $E=200$\,GPa; bays of 6.0, 7.5 or 9.0\,m; a taller first storey (4.0--4.5\,m) over typical storeys (3.2--3.8\,m); gravity loads from $1.2D+1.0L$ area loads ($D=3.5$--$5.0$\,kPa, $L=2.0$--$3.0$\,kPa, lighter roof) on a 4--6\,m tributary width; wind forces from a uniform facade pressure of 0.6--1.2\,kPa. Beams are sized for an $L/360$ live-load deflection and a simple-span moment demand, columns for an $h/400$ storey drift under wind (shear-building estimate with a beam-flexibility factor of 0.5) with one section per two storeys, and 80\% of the bases are fixed. Section properties are clipped to the training ranges, so only the parameter correlations and load proportions are new; line loads reach 52\,kN/m, above the training range, which the load scaling of Section~\ref{sec:scaling} makes irrelevant.
\end{itemize}
The OOD-G set tests extrapolation in topology and size, the OOD-L set extrapolation in the load pattern, and the DES set a realistic, design-driven corner of the design space; none is used for training, model selection or early stopping. The sampling procedure is listed step by step in \ref{app:sampling}.

\begin{table}[t]
\centering
\caption{Sampling ranges of the frame family. $\mathcal{U}$: uniform, $\log\mathcal{U}$: log-uniform. Units: kN, m.}\label{tab:ranges}
\footnotesize
\begin{adjustbox}{max width=\textwidth}
\begin{tabular}{p{3.4cm}p{5.4cm}p{4.6cm}}
\toprule
Parameter & Training design space (ID) & Unseen scenarios \\
\midrule
Storeys $n_s$ & $\{1,\dots,5\}$ & OOD-G: $\{6,7,8\}$ or $n_b\in\{5,6\}$ \\
Bays $n_b$ & $\{1,\dots,4\}$ & OOD-G: up to $6$ \\
Bay width $L_b$ & $\mathcal{U}(3,9)$ & same \\
Storey height $h_s$ & $\mathcal{U}(2.8,4.5)$ & OOD-G: first storey $\mathcal{U}(4.5,6.0)$ in 50\% \\
Young's modulus $E$ (kN/m$^2$) & $\log\mathcal{U}(2.5\times10^7, 2.1\times10^8)$ & same \\
Column $I_{c}$ at base (m$^4$) & $\log\mathcal{U}(5\times10^{-5},2\times10^{-3})$, taper $\mathcal{U}(0.5,1)$ & same \\
Beam $I_{b,s}$ (m$^4$) & $\log\mathcal{U}(5\times10^{-5},2\times10^{-3})$ per storey & same \\
Radius of gyration $r$ & $\mathcal{U}(0.08,0.30)$ (columns, beams) & same \\
Base condition & fixed (70\%) / pinned (30\%) & same \\
Gravity load $w$ (kN/m) & $\mathcal{U}(5,40)$ per floor, all bays & OOD-L: 30--60\% of $w$ in random bays \\
Lateral profile & uniform or inverted triangle & OOD-L: parabolic or roof-only \\
Roof lateral force (kN) & $\mathcal{U}(5,80)$, random sign & same \\
\bottomrule
\end{tabular}
\end{adjustbox}
\end{table}

\subsection{PyNite analyses and statics verification}
Each frame is built as an \texttt{FEModel3D} in PyNite~\cite{pynite} (version~3.2.0 of the fork cited) with one Euler--Bernoulli member per column and beam segment. Out-of-plane translations and rotations are restrained at every node so that the 3D solver reproduces planar behaviour. The lateral and gravity loads are defined as two separate load cases, W and G, and solved by a single first-order linear analysis (\texttt{analyze\_linear}). For each case the nodal displacements $(u_x,u_y,\theta_z)$ and the global member end forces are extracted; the latter are converted to the basic forces of Section~\ref{sec:statics}. The wall-clock time of the complete PyNite call chain (model construction, analysis and result extraction) is recorded for every frame and later used for the cost comparison.

As a guard against convention errors, every FE result is checked against an independent implementation of the frame statics (Section~\ref{sec:statics}): the out-of-balance nodal force computed from the PyNite member forces must vanish. All 7{,}500 analyses pass with residuals at round-off level (below $10^{-6}$ of the largest applied load; typically $10^{-11}$\,kN). The data comprise a candidate pool of 5{,}000 ID frames, 400 ID validation frames, three test sets of 600 frames each (ID, OOD-G, OOD-L) and the 300 DES frames, generated with disjoint random seeds.

\subsection{Independent verification of the reference solutions}\label{sec:validation}
Because every error reported below is measured against PyNite, the reference data were verified independently in two ways. First, all 2{,}100 test frames were re-analysed with two solvers that share no code with PyNite: a direct-stiffness solver for planar Euler--Bernoulli frames written for this study (vectorised element matrices, consistent fixed-end forces for the member loads, sparse direct solution with SciPy), and OpenSeesPy~\cite{zhu2018openseespy} (version 3.8.0; \texttt{elasticBeamColumn} elements with a linear transformation, \texttt{beamUniform} element loads, UMFPACK solver). Table~\ref{tab:validation} lists the largest relative $L_2$ difference per frame between each solver and PyNite for horizontal and vertical displacements, rotations, support reactions, column and beam end moments, end shears and axial forces. All differences are below $1.3\times10^{-11}$, i.e.\ at round-off level. Second, PyNite was run on three classical cases with closed-form solutions: a portal frame with a rigid beam under a lateral load $P$ with fixed bases (sway $Ph^3/(24EI_c)$, column end moments $Ph/4$) and with pinned bases (sway $Ph^3/(6EI_c)$, top moments $Ph/2$), and a beam with fixed ends under a uniform load (end moments $wL^2/12$). The PyNite results differ from the closed-form values by at most $9\times10^{-6}$ in relative terms, which is explained by the finite stiffness of the very stiff members used to represent the rigid idealisations.

\begin{table}[t]
\centering
\caption{Independent verification of the PyNite reference solutions. Top: largest relative $L_2$ difference per frame between each solver and PyNite over every frame of each test set. Bottom: closed-form benchmarks analysed with PyNite.}\label{tab:validation}
\footnotesize
\begin{adjustbox}{max width=\textwidth}
\inputtable{tables/tab_validation.tex}
\end{adjustbox}
\end{table}

%% ---------------------------------------------------------------------------
\section{Mechanics-constrained surrogate}\label{sec:method}

The surrogate maps the frame description to the full response field: all nodal displacements and the basic forces of all members. Three kinds of mechanics are built into it, in decreasing order of strictness: (i) exact constraints that hold for any network output (boundary conditions, homogeneity in the load intensities, nodal equilibrium after projection); (ii) mechanics-derived inputs that are exact statical quantities or classical approximations; and (iii) an inductive bias in the architecture (message passing on the structural graph).

\subsection{Statics of the frame and the equilibrium constraint}\label{sec:statics}
For member $e$ between nodes $i$ and $j$ with length $L_e$ and direction cosines $(c_e,s_e)$, the basic forces are $\mathbf{q}_e = [N_e, M_{i,e}, M_{j,e}]^{\mathsf T}$: the axial force (tension positive) and the two end moments acting on the member, counter-clockwise positive. Columns run from bottom to top and beams from left to right; nodal displacements $(u_x,u_y,\theta_z)$ are positive along global $X$, $Y$ (upwards) and counter-clockwise (Fig.~\ref{fig:signs}). For a member carrying a transverse distributed load $w_e$, the six local end forces are
\begin{equation}
\mathbf{f}_e = \mathbf{b}_e\mathbf{q}_e + \mathbf{f}^{0}_e, \qquad
\mathbf{b}_e = \begin{bmatrix} -1 & 0 & 0\\ 0 & 1/L_e & 1/L_e \\ 0 & 1 & 0 \\ 1 & 0 & 0 \\ 0 & -1/L_e & -1/L_e \\ 0 & 0 & 1 \end{bmatrix},
\end{equation}
where $w_e$ is the local transverse load intensity (positive along the local $y$ axis, so gravity on a beam is negative) and $\mathbf{f}^{0}_e = -\tfrac{1}{2}w_eL_e\,[0,1,0,0,1,0]^{\mathsf T}$ is the simply supported particular solution.

\begin{figure}[t]
\centering
\begin{tikzpicture}[>=Stealth, font=\footnotesize, scale=1.0]
% global axes
\draw[->] (-0.6,-0.4) -- (0.2,-0.4) node[right]{$X$};
\draw[->] (-0.6,-0.4) -- (-0.6,0.4) node[above]{$Y$};
% member
\coordinate (I) at (1.2,0.6); \coordinate (J) at (6.2,0.6);
\draw[very thick] (I) -- (J);
\fill (I) circle (1.6pt) node[below left]{$i$}; \fill (J) circle (1.6pt) node[below right]{$j$};
\draw[->, gray] (3.2,0.25) -- (4.2,0.25) node[right]{local $x$};
\draw[->, gray] (3.2,0.25) -- (3.2,-0.35) node[below]{};
\draw[->, gray] (3.2,0.25) -- (3.2,0.0);
% distributed load
\foreach \x in {1.7,2.4,...,5.9} \draw[->] (\x,1.35) -- (\x,0.7);
\draw (1.4,1.35) -- (6.0,1.35); \node at (3.7,1.6) {$w_e<0$ (gravity, local $y$ up)};
% axial forces (tension positive)
\draw[->, thick, blue!70!black] (I) -- ++(-0.8,0) node[left]{$N_e$};
\draw[->, thick, blue!70!black] (J) -- ++(0.8,0) node[right]{$N_e$};
% end moments CCW
\draw[->, thick, red!70!black] ([shift={(0.45,0)}]I) arc (0:250:0.45) node[pos=0.35, above left]{$M_{i,e}$};
\draw[->, thick, red!70!black] ([shift={(0.45,0)}]J) arc (0:250:0.45) node[pos=0.25, above right]{$M_{j,e}$};
% node DOFs
\begin{scope}[shift={(8.4,0.6)}]
\fill (0,0) circle (1.6pt);
\draw[->, thick] (0,0) -- (0.9,0) node[right]{$u_x$};
\draw[->, thick] (0,0) -- (0,0.9) node[above]{$u_y$};
\draw[->, thick] (0.45,0) arc (0:270:0.45) node[pos=0.55, left]{$\theta_z$};
\node[below] at (0,-0.55) {nodal DOFs};
\end{scope}
\end{tikzpicture}
\caption{Sign conventions. Basic forces of member $e$ ($i\to j$): axial force $N_e$, tension positive; end moments $M_{i,e}$, $M_{j,e}$ acting on the member, counter-clockwise positive; transverse load $w_e$ along the local $y$ axis. Nodal displacements $u_x$, $u_y$ along the global axes and rotation $\theta_z$ counter-clockwise. Columns are oriented bottom to top, beams left to right.}\label{fig:signs}
\end{figure} This parametrisation satisfies member equilibrium identically \cite{mcguire2000matrix}. Rotating to global axes and assembling at the nodes gives the nodal equilibrium equations at the $n_f$ free degrees of freedom (DOFs),
\begin{equation}\label{eq:equilibrium}
\mathbf{C}\,\mathbf{q} = \mathbf{d}, \qquad \mathbf{C}\in\mathbb{R}^{n_f\times 3m},\quad \mathbf{d} = \mathbf{P}_f - \mathbf{p}^0_f ,
\end{equation}
where $\mathbf{q}$ stacks the basic forces of all $m$ members, $\mathbf{P}_f$ are the applied nodal loads and $\mathbf{p}^0_f$ the assembled particular solutions. $\mathbf{C}$ is the equilibrium (statics) matrix. For a kinematically stable frame it has full row rank and $3m-n_f$ equals the degree of static indeterminacy, so Eq.~\eqref{eq:equilibrium} defines an affine set of dimension $3m-n_f>0$ that contains the exact FE solution but also every other statically admissible force field. Equilibrium alone therefore cannot determine the forces, which is precisely why a learned prior is needed; conversely, any force prediction that violates Eq.~\eqref{eq:equilibrium} is mechanically inadmissible. Support reactions follow from $\mathbf{q}$ at the restrained DOFs, so a force field satisfying Eq.~\eqref{eq:equilibrium} also satisfies global equilibrium. Assembling $\mathbf{C}$ costs $\mathcal{O}(m)$ and requires no stiffness matrix.

\subsection{Dimensional analysis and exact load-case superposition}\label{sec:scaling}
For a linear elastic frame the response is linear in the loads. If $E$ is the same in every member, as here, the displacements and rotations are inversely proportional to $E$, whereas the member forces of this statically indeterminate structure depend only on the relative stiffnesses and are therefore independent of a uniform $E$. The surrogate imposes these properties exactly through dimensional analysis~\cite{buckingham1914pi}. Let $H=\frac{1}{n_s}\sum_s h_s$ be the mean storey height and $I_\mathrm{ref}=\frac{1}{n_s}\sum_s I_{c,s}$ the mean column second moment of area, so that $EI_\mathrm{ref}$ is a reference flexural rigidity. Each load case $c\in\{\mathrm{W},\mathrm{G}\}$ is scaled by its own intensity: the signed base shear $\lambda_\mathrm{W}=\sum_s F_s$ and the total gravity load $\lambda_\mathrm{G}=\sum_e |w_e|L_e$. The response is written as
\begin{equation}\label{eq:superposition}
\begin{aligned}
(u_x,u_y) &= \sum_{c} \frac{\lambda_c H^3}{EI_\mathrm{ref}}\,(\hat u_{x,c},\hat u_{y,c})(\boldsymbol{\pi}_c), &
\theta_z &= \sum_{c} \frac{\lambda_c H^2}{EI_\mathrm{ref}}\,\hat\theta_{c}(\boldsymbol{\pi}_c),\\
N &= \sum_c \lambda_c\, \hat N_c(\boldsymbol{\pi}_c), &
M &= \sum_c \lambda_c H\,\hat M_c(\boldsymbol{\pi}_c),
\end{aligned}
\end{equation}
where hats denote the network outputs and $\boldsymbol{\pi}_c$ is a set of dimensionless input groups (Table~\ref{tab:dimensions}): lengths divided by $H$, second moments of area divided by $I_\mathrm{ref}$, areas as $AH^2/I_\mathrm{ref}$, lateral loads divided by $\lambda_\mathrm{W}$ and gravity loads by $\lambda_\mathrm{G}$. Because $\boldsymbol{\pi}_c$ is invariant to the load magnitudes and to $E$, Eq.~\eqref{eq:superposition} is exactly homogeneous of degree one in each load-case intensity, the displacements scale exactly with $1/E$, the forces are exactly independent of a uniform $E$, and reversing the lateral load reverses the lateral response. $E$ does not appear in the force expressions at all, and $\lambda_c$ cancels in every input group.

In the generator, $|\lambda_\mathrm{W}|\ge 5$\,kN (all floor forces of a frame have the same sign) and $\lambda_\mathrm{G}>0$, so neither scale can vanish. In deployment a load case with $\lambda_c=0$ is simply omitted from the sum; a lateral pattern whose floor forces cancel ($\sum_sF_s\approx0$ but $\sum_s|F_s|>0$) would make $\lambda_\mathrm{W}$ ill-conditioned and should be scaled by $\sum_s|F_s|$ instead, a case that does not occur in this study and that the domain check of Section~\ref{sec:ood} would flag as unseen. Predicting the two cases separately also avoids a scaling problem: in the combined case the lateral sway is often orders of magnitude smaller than the gravity-induced quantities, so a single scale under-weights it. The ablation in Section~\ref{sec:ablation} quantifies this choice. Restrained DOFs are multiplied by zero, so the boundary conditions also hold exactly.

\begin{table}[t]
\centering
\caption{Dimensional consistency of the scaled formulation (units kN, m). Every input and output group of the network is dimensionless and independent of $\lambda_c$ and $E$.}\label{tab:dimensions}
\footnotesize
\begin{tabular}{llll}
\toprule
Quantity & Units & Scale & Dimensionless group \\
\midrule
Lengths $x$, $L$, $h$ & m & $H$ & $x/H$, $L/H$, $h/H$ \\
Second moment of area $I$ & m$^4$ & $I_\mathrm{ref}$ & $I/I_\mathrm{ref}$ \\
Cross-sectional area $A$ & m$^2$ & $I_\mathrm{ref}/H^2$ & $AH^2/I_\mathrm{ref}$ \\
Floor force $F_s$, storey shear $V_s$ & kN & $\lambda_\mathrm{W}$ & $F_s/\lambda_\mathrm{W}$, $V_s/\lambda_\mathrm{W}$ \\
Overturning moment & kN\,m & $\lambda_\mathrm{W}H$ & $\mathrm{OTM}/(\lambda_\mathrm{W}H)$ \\
Line load $w$ (resultant $wL$) & kN/m (kN) & $\lambda_\mathrm{G}$ & $wL/\lambda_\mathrm{G}$ \\
Drift estimate $\delta_s$ & m & $\lambda_\mathrm{W}H^3/EI_\mathrm{ref}$ & $\delta_s EI_\mathrm{ref}/(\lambda_\mathrm{W}H^3)$ \\
Displacements $u_x$, $u_y$ & m & $\lambda_cH^3/EI_\mathrm{ref}$ & $\hat u_{x,c}$, $\hat u_{y,c}$ \\
Rotation $\theta_z$ & rad & $\lambda_cH^2/EI_\mathrm{ref}$ & $\hat\theta_c$ \\
Axial force $N$ & kN & $\lambda_c$ & $\hat N_c$ \\
End moment $M$ & kN\,m & $\lambda_cH$ & $\hat M_c$ \\
Equilibrium rows (force / moment) & kN / kN\,m & $\lambda_c$ / $\lambda_cH$ & rows of $\mathbf{C}$, $\mathbf{d}_c$ \\
\bottomrule
\end{tabular}
\end{table}

\subsection{Mechanics-informed features}\label{sec:features}
Each node carries 25 dimensionless features and each member 10. Beyond position, topology and section ratios, they include two quantities from elementary statics that are exact and cost $\mathcal{O}(n)$, the storey shear $V_s=\sum_{k\ge s}F_k$ and the overturning moment of the loads above the storey, and three classical approximations. The first is the gravity load carried by each column line computed from tributary areas, i.e.\ as if every beam were simply supported; in a rigid frame continuity and sway redistribute these loads, so the feature is an approximation, not an exact statical quantity. The other two are the beam-to-column stiffness ratio of each storey, $\rho_s = (n_b I_{b,s}/\bar L_b)/((n_b+1)I_{c,s}/h_s)$, and the shear-building drift estimate $\delta_s = V_s/k_s$ with $k_s=(n_b+1)\,12EI_{c,s}/h_s^3$ (or $3EI/h^3$ for a pinned first storey), together with its height-wise cumulative sum. The network learns the correction from these simple estimates to the FE response, which is a smoother target than the response itself. The full feature dictionary with formulas is given in \ref{app:features}. The same features are given to all tabular baselines.

\subsection{Graph network}\label{sec:gnn}
Frames of different size and topology are represented as graphs whose nodes are joints and whose edges are members, and processed with an encode--process--decode message passing network~\cite{battaglia2018relational,gilmer2017mpnn,sanchez2020learning,pfaff2021meshgraphnets}. Node and edge features are encoded by two-layer perceptrons (width 96, SiLU, layer normalisation). Each of the eight processing steps updates every edge from its own state and the states of its end nodes, $\mathbf{z}_e \leftarrow \mathbf{z}_e + \phi_e([\mathbf{z}_e,\mathbf{h}_i,\mathbf{h}_j])$, and every node from the sums of the messages of the members that start and end at it, kept separate so that the network can distinguish the two ends of a member, $\mathbf{h}_i \leftarrow \mathbf{h}_i + \phi_n([\mathbf{h}_i,\sum_{e:i=\mathrm{start}}\mathbf{z}_e,\sum_{e:i=\mathrm{end}}\mathbf{z}_e])$. Eight steps nearly cover the graph diameter of the training frames (at most nine members), and frame-level information is supplied in addition through the statics features; the effect of the depth (4, 8, 12 and 16 steps) is examined in Section~\ref{sec:graphbaselines}. Decoders output six dimensionless displacements per node and six dimensionless basic forces per member (three per load case). The loss is the mean squared error of the outputs standardised per component, over free DOFs only. Training uses AdamW~\cite{loshchilov2019adamw} with a one-cycle learning-rate schedule (peak $2\times10^{-3}$), batches of 32 frames, gradient clipping at 1.0 and 150 epochs, implemented in PyTorch~\cite{paszke2019pytorch}. A variant with a soft equilibrium penalty $\|\mathbf{C}\hat{\mathbf{q}}_c-\mathbf{d}_c\|^2$ added to the loss is examined in the ablation.

\subsection{Equilibrium projection}\label{sec:projection}
A data-driven prediction $\hat{\mathbf{q}}$ will in general violate Eq.~\eqref{eq:equilibrium}. The proposed constraint layer replaces it with the closest statically admissible force field in a metric $\boldsymbol{\Sigma}^{-1}$:
\begin{equation}\label{eq:projection}
\mathbf{q}^\star = \arg\min_{\mathbf{C}\mathbf{q}=\mathbf{d}} (\mathbf{q}-\hat{\mathbf{q}})^{\mathsf T}\boldsymbol{\Sigma}^{-1}(\mathbf{q}-\hat{\mathbf{q}})
= \hat{\mathbf{q}} + \boldsymbol{\Sigma}\mathbf{C}^{\mathsf T}\left(\mathbf{C}\boldsymbol{\Sigma}\mathbf{C}^{\mathsf T}\right)^{-1}\left(\mathbf{d}-\mathbf{C}\hat{\mathbf{q}}\right),
\end{equation}
which follows from the Karush--Kuhn--Tucker conditions of the equality-constrained quadratic programme~\cite{nocedal2006numerical}. It is applied separately to each load case in dimensionless form. Two metrics are compared: the Euclidean metric on the dimensionless forces ($\boldsymbol{\Sigma}=\mathbf{I}$), and the \emph{variance-weighted} metric built from the ensemble variance $\sigma_k^2$ of each predicted dimensionless force component $k$ (Section~\ref{sec:uq}), so that the components the model is least sure of absorb most of the correction. Because an ensemble variance can be arbitrarily close to zero, the metric is regularised as
\begin{equation}\label{eq:sigmareg}
\boldsymbol{\Sigma}_\mathrm{reg} = \mathrm{diag}\!\left(\sigma_k^2 + \varepsilon\,\bar{\sigma}^2\right), \qquad \bar\sigma^2 = \frac{1}{3m}\sum_k \sigma_k^2,\quad \varepsilon = 10^{-6},
\end{equation}
per load case and frame; the floor makes $\boldsymbol{\Sigma}_\mathrm{reg}$ positive definite and scale-free. The system $(\mathbf{C}\boldsymbol{\Sigma}\mathbf{C}^{\mathsf T})\boldsymbol{\mu} = \mathbf{d}-\mathbf{C}\hat{\mathbf{q}}$ is solved by a dense LU factorisation with partial pivoting (LAPACK \texttt{gesv} through NumPy), so there is no iterative tolerance; the resulting equilibrium residuals are at round-off level (Section~\ref{sec:results}). The sensitivity of the results to $\varepsilon$, to the alternative floor $\mathrm{diag}(\max(\sigma_k^2,\varepsilon\bar\sigma^2))$ and the condition numbers of $\mathbf{C}\boldsymbol{\Sigma}\mathbf{C}^{\mathsf T}$ are reported in Section~\ref{sec:sensitivity}. If the prediction error were Gaussian, unbiased and independent between components with covariance $\boldsymbol{\Sigma}$, Eq.~\eqref{eq:projection} would be the posterior mean after conditioning on the equilibrium equations. The ensemble variance is at best a crude diagonal approximation of that covariance (the errors of neighbouring members are correlated and the ensemble under-represents the error under distribution shift), so this interpretation motivates the choice of metric but does not make $\mathbf{q}^\star$ a calibrated posterior.

\begin{proposition}\label{prop:projection}
Let $\mathbf{q}_\mathrm{FE}$ be the exact solution, so that $\mathbf{C}\mathbf{q}_\mathrm{FE}=\mathbf{d}$. For any prediction $\hat{\mathbf{q}}$ and any symmetric positive definite $\boldsymbol{\Sigma}$, the projected forces satisfy
$\|\mathbf{q}^\star-\mathbf{q}_\mathrm{FE}\|^2_{\boldsymbol{\Sigma}^{-1}} = \|\hat{\mathbf{q}}-\mathbf{q}_\mathrm{FE}\|^2_{\boldsymbol{\Sigma}^{-1}} - \|\hat{\mathbf{q}}-\mathbf{q}^\star\|^2_{\boldsymbol{\Sigma}^{-1}} \le \|\hat{\mathbf{q}}-\mathbf{q}_\mathrm{FE}\|^2_{\boldsymbol{\Sigma}^{-1}}$.
\end{proposition}
\begin{proof}
$\mathbf{q}^\star$ is the $\boldsymbol{\Sigma}^{-1}$-orthogonal projection of $\hat{\mathbf{q}}$ onto the affine set $\mathcal{A}=\{\mathbf{q}:\mathbf{C}\mathbf{q}=\mathbf{d}\}$. Since $\mathbf{q}^\star-\mathbf{q}_\mathrm{FE}$ lies in $\ker\mathbf{C}$ and $\hat{\mathbf{q}}-\mathbf{q}^\star=\boldsymbol{\Sigma}\mathbf{C}^{\mathsf T}\boldsymbol{\mu}$ for some $\boldsymbol{\mu}$, their $\boldsymbol{\Sigma}^{-1}$-inner product is $\boldsymbol{\mu}^{\mathsf T}\mathbf{C}(\mathbf{q}^\star-\mathbf{q}_\mathrm{FE})=0$, and the identity follows from Pythagoras' theorem.
\end{proof}
The projection can therefore never increase the error measured in its own metric, and it removes exactly the part of the error that is visible to statics. It does not guarantee a reduction of every individual component (for example the end moments alone), which is why the choice of metric matters and is examined empirically. The projection needs one dense solve of size $n_f$ per load case, which for the frames considered here takes a fraction of a millisecond, and it is model-agnostic: it is applied below to the baselines as well.

\subsection{Displacement--force compatibility}\label{sec:compat}
The network predicts displacements and forces with separate decoders, and the projection changes only the forces. The two fields are therefore not linked by the element stiffness relations. To measure this, the basic forces implied by the predicted displacements are computed member by member as $\mathbf{q}_\mathrm{const}(\hat{\mathbf{u}})$, i.e.\ the local end forces $\mathbf{k}_e\mathbf{T}_e\hat{\mathbf{u}}_e + \mathbf{f}^\mathrm{fix}_e$ (element stiffness $\mathbf{k}_e$, rotation $\mathbf{T}_e$ and fixed-end forces $\mathbf{f}^\mathrm{fix}_e$ of the member load), and the compatibility residual is defined as
\begin{equation}\label{eq:compat}
r_\mathrm{compat} = \frac{\|\mathbf{S}(\mathbf{q}-\mathbf{q}_\mathrm{const}(\hat{\mathbf{u}}))\|_2}{\|\mathbf{S}\,\mathbf{q}_\mathrm{const}(\hat{\mathbf{u}})\|_2},
\qquad r_{\mathrm{compat},M} = \frac{\|\mathbf{M}-\mathbf{M}_\mathrm{const}(\hat{\mathbf{u}})\|_2}{\|\mathbf{M}_\mathrm{const}(\hat{\mathbf{u}})\|_2},
\end{equation}
where $\mathbf{q}$ is the force field being assessed (raw, projected or recovered) and $\mathbf{S}$ divides the moment entries by $H$ so that forces and moments are commensurate. $r_{\mathrm{compat},M}$ uses the end moments only, because the axial forces implied by displacements, $EA/L$ times the member elongation, amplify small displacement errors by the large axial stiffness. A zero residual means that the forces are those the predicted displacements would produce; the exact FE solution has $r_\mathrm{compat}=0$ and zero equilibrium residual simultaneously. Recovering the forces from the predicted displacements, as in displacement-based surrogates~\cite{leng2026mcgnn}, gives $r_\mathrm{compat}=0$ by construction but sacrifices equilibrium; projecting the directly predicted forces gives exact equilibrium but leaves compatibility unenforced. Both options, and the effect of the projection on $r_\mathrm{compat}$, are quantified in Section~\ref{sec:compatresults}. A field that satisfies both conditions exactly is the FE solution itself; the predicted displacements can then serve as the initial iterate of an iterative solver, which is also examined.

\subsection{Predictive uncertainty}\label{sec:uq}
Epistemic uncertainty~\cite{derkiureghian2009aleatory} is estimated with a deep ensemble~\cite{lakshminarayanan2017ensembles} of $M=5$ networks that differ in initialisation and batch order and share the training data and normalisation. The ensemble mean is the prediction and the ensemble variance $\sigma^2$ of each output feeds the projection metric. For decisions at the level of a whole frame, the relative spread of the end moments,
\begin{equation}\label{eq:score}
s = \frac{\|\boldsymbol{\sigma}_M\|_2}{\|\hat{\mathbf{M}}\|_2},
\end{equation}
is used, where $\boldsymbol{\sigma}_M$ collects the physical standard deviations of all end moments. A second, label-free indicator is also examined: the equilibrium residual $\|\mathbf{C}\hat{\mathbf{q}}-\mathbf{d}\|/\|\mathbf{d}\|$ of the unprojected ensemble mean, which measures how far the raw prediction is from being statically admissible. The two indicators have different roles in the workflow. The ensemble spread $s$ is used at deployment, to decide whether an individual prediction is accepted or referred to the FE solver; the equilibrium residual is used during training, to decide which candidate frames are worth an FE analysis (Section~\ref{sec:workflow}).

\subsection{Domain check and calibrated referral threshold}\label{sec:ood}
An ensemble can be confidently wrong on inputs unlike its training data~\cite{ovadia2019shift}, and Section~\ref{sec:uqresults} shows that this happens for the unseen load patterns. The deployment rule therefore has two stages, both calibrated on the 400 ID validation frames, which are used neither for training nor for testing.

\emph{Stage 1: domain check.} Each frame is summarised by ten geometry descriptors (numbers of storeys and bays, largest storey height, first-to-mean storey height ratio, mean bay width, ratio of the largest to the smallest bay, height-to-width aspect ratio, $\log_{10}E$, $\log_{10}$ of the beam-to-column second-moment ratio, base fixity) and five load-pattern descriptors (normalised height of the centroid of the lateral forces, roof share of the lateral load, mean ratio of each beam load to the largest load on its floor, mean within-floor coefficient of variation of the beam loads, and $\log_{10}$ of the gravity-to-lateral load ratio). All are available before any analysis. The descriptors are standardised with the training-set statistics and the Mahalanobis distance~\cite{lee2018mahalanobis} to the training distribution is computed with a Ledoit--Wolf shrinkage covariance~\cite{ledoit2004shrinkage} estimated on the 3{,}000 training frames. A frame is outside the domain, and referred to FE, if its distance exceeds the 99th percentile of the distances of the validation frames (a nominal false-referral rate of 1\% for in-distribution frames). Distances based on the geometry and the load descriptors alone are also reported, to show which group detects which shift.

\emph{Stage 2: conformal error bound.} For frames inside the domain, the ensemble spread is converted into an error bound by split conformal prediction~\cite{lei2018conformal}. On the in-domain validation frames, the non-conformity scores $a_i = e_{M,i}/s_i$ (true moment error over spread) are computed and $\hat a$ is their $\lceil (n+1)(1-\alpha)\rceil$-th smallest value with $\alpha=0.1$. For a new frame the bound is $\hat a\,s$, which under exchangeability with the calibration frames covers the true moment error with probability at least $1-\alpha$. A frame is accepted if $\hat a\,s\le\tau$ for a tolerance $\tau$ on the relative moment error (5\% and 10\% are reported) and referred otherwise. Unlike the risk--coverage curves of Section~\ref{sec:uqresults}, which are computed on the test sets and serve only to describe the ranking quality, both thresholds of this rule are fixed before the test data are seen.

%% ---------------------------------------------------------------------------
\section{Uncertainty-guided workflow}\label{sec:workflow}

The surrogate is embedded in a workflow with two uses of the uncertainty estimate (Fig.~\ref{fig:workflow}).

\paragraph{Adaptive FE data generation} Training starts from $N_0$ randomly chosen frames of the candidate pool, analysed with PyNite. In each round an ensemble is trained, every remaining candidate is scored without any FE analysis, the $B$ highest-scoring candidates are analysed with PyNite and added to the training set. This is pool-based active learning~\cite{settles2009active} in the spirit of adaptive surrogate methods in structural reliability~\cite{bichon2008egra,echard2011akmcs}. Four acquisition rules are compared: random sampling, the ensemble variance (mean normalised variance over all outputs), the equilibrium residual of the raw prediction, and a diversity-aware variant of the latter. Selecting the $B$ largest scores can pick many near-identical frames from one corner of the design space, so the diversity-aware rule shortlists the $3B$ candidates with the largest residuals, groups them into $B$ clusters by $k$-means on standardised frame descriptors (the mean and maximum of the node features) and takes the highest-residual frame of each cluster, in the spirit of diverse mini-batch active learning~\cite{zhdanov2019diverse,sener2018coreset}. Each rule is repeated with five random initial sets (seeds 0--4). The candidates are labelled in advance for convenience, but a label is revealed to the learner only when the strategy requests that analysis, and the recorded PyNite time of each requested analysis is charged to the budget.

\paragraph{FE referral at deployment} For a batch of new frames, all are first predicted and projected. A frame is sent to PyNite, and its surrogate result discarded, if it fails the domain check or if its conformal error bound exceeds the tolerance (Section~\ref{sec:ood}); the remaining predictions are accepted. To describe how well the spread ranks the errors independently of any threshold, risk--coverage curves~\cite{geifman2017selective} are also reported, in which the fraction of frames with the largest spread referred to FE is varied.

\begin{figure}[t]
\centering
\begin{tikzpicture}[node distance=4mm and 11mm, every node/.style={font=\footnotesize, align=center},
  box/.style={draw, rounded corners=2pt, minimum height=9mm, text width=18mm, fill=blue!4},
  fe/.style={box, fill=orange!10}, arr/.style={-{Stealth[length=2mm]}, thick}]
\node[box] (param) {Frame parameters\\(geometry, sections,\\loads, supports)};
\node[box, right=of param] (feat) {Dimensionless\\graph + statics\\features};
\node[box, right=of feat] (gnn) {GNN ensemble\\(per load case)};
\node[box, right=of gnn] (proj) {Superposition +\\equilibrium\\projection};
\node[box, below=9mm of proj] (score) {Domain check +\\conformal bound $\hat a s$};
\node[box, below=9mm of gnn, fill=green!8] (accept) {Accept surrogate\\response};
\node[fe, below=9mm of feat] (fe) {PyNite FE\\analysis};
\node[fe, below=9mm of param] (train) {Training set\\(adaptive growth)};
\draw[arr] (param) -- (feat); \draw[arr] (feat) -- (gnn); \draw[arr] (gnn) -- (proj);
\draw[arr] (proj) -- (score);
\draw[arr] (score) -- node[above, font=\scriptsize]{accept} (accept);
\draw[arr] (score.south) -- ++(0,-4mm) -| node[pos=0.25, below]{out of domain or $\hat a s>\tau$: refer} (fe.south);
\draw[arr] (fe) -- (train);
\draw[arr] (train.north) -- node[left, align=right]{retrain;\\acquire by\\eq.\ residual} (param.south);
\end{tikzpicture}
\caption{Proposed workflow. Exact mechanics (superposition, boundary conditions, static equilibrium) is imposed around a learned graph ensemble. During training, new FE analyses are chosen by the equilibrium residual of the raw prediction; at deployment, a domain check and a conformal bound on the ensemble spread decide which frames are referred to PyNite.}\label{fig:workflow}
\end{figure}

%% ---------------------------------------------------------------------------
\section{Baselines and evaluation protocol}\label{sec:protocol}

\paragraph{Baselines} Five baselines represent conventional regression and widely used machine learning: ridge regression~\cite{hoerl1970ridge} on the features (linear) and on all their pairwise products (quadratic), random forests~\cite{breiman2001rf} (200 trees), histogram gradient boosting~\cite{friedman2001gbm} (600 iterations, one model per output) and a three-layer perceptron (256 units per layer, early stopping), all from scikit-learn~\cite{pedregosa2011sklearn}. Because fixed-length regressors cannot ingest graphs of varying size, they predict per node (displacements, from the 25 node features) and per member (basic forces, from the 10 member features concatenated with the features of both end nodes). They receive the same dimensionless, superposition-split targets and the same hard boundary conditions as the graph network, so that the comparison isolates the learning model and the constraint layer rather than the preprocessing. The quadratic ridge model is fitted on a random subset of 30{,}000 rows to keep the design matrix in memory. All models are trained on the same first 3{,}000 frames of the ID pool.

\paragraph{Graph-learning baselines}\label{sec:graphbaselinesdef} Five graph models isolate the design choices of the proposed network. All are trained on the same 3{,}000 frames for 150 epochs with three random seeds each.
(i) A \emph{conventional GNN}: the same architecture, but with physical inputs (coordinates, support flags, nodal loads, member length, $\log_{10}E$, $\log_{10}I$, $\log_{10}A$, line load) and physical targets for the combined load case, standardised over the training set; no statics features, no load-case split and no dimensional scaling.
(ii) A \emph{global-pooling} variant of the proposed network, in which the graph-wide mean of the node states is appended to every node update, giving every node frame-level context at every step.
(iii) A \emph{StructGNN-style} network following Chou et al.~\cite{chou2024structgnn}: one pseudo node per storey linked to all joints of that floor and to the pseudo node below, linear encoders to width 256, linear messages with mean aggregation in both edge directions, $n_s+1$ propagation steps for a frame with $n_s$ storeys (mini-batches grouped by $n_s$), decoders with one ReLU hidden layer and a mean-absolute-error loss. To allow a fair comparison it receives the same dimensionless features and targets as the proposed network. Two deviations from the original are needed: the weights are shared across propagation steps, because the unseen-geometry frames need up to nine steps whereas training frames need at most six, and a residual connection is added to train stably with the learning-rate schedule used for all models (the original used 1{,}000 epochs with batch size one, which was not affordable for all seeds here).
(iv) A \emph{displacement-only} network: the proposed architecture trained on the displacements only, with member forces recovered from the predicted displacements through the element stiffness relations and the fixed-end forces.
(v) A \emph{Leng-style} network following Leng et al.~\cite{leng2026mcgnn}: as (iv), with the normalised free-DOF stiffness residual $\|\mathbf{K}_{ff}\hat{\mathbf{u}}_f-\mathbf{f}_f\|^2/\|\mathbf{f}_f\|^2$ added to the loss with weight $10^{-2}$, chosen from $\{10^{-4},10^{-3},10^{-2},10^{-1},1\}$ on a short pilot run because larger weights let the poorly conditioned axial rows dominate. These are re-implementations from the published descriptions, not the authors' code.
The soft-equilibrium-penalty and no-superposition ablations of the proposed network and a depth ablation (4, 8, 12, 16 message-passing steps) are trained with three seeds as well.

\paragraph{Metrics} All errors are computed per frame in physical units against PyNite and summarised over a test set by the median and the 95th percentile (p95). With $\hat{\cdot}$ the surrogate and $\|\cdot\|$ the Euclidean norm over all entries of a frame, the field errors are relative $L_2$ errors, $e_a = \|\hat{\mathbf{a}}-\mathbf{a}\|/\|\mathbf{a}\|$, for the horizontal displacements ($e_{u_x}$), vertical displacements ($e_{u_y}$) and rotations ($e_\theta$) of all nodes, the end moments ($e_M$), axial forces ($e_N$) and end shears ($e_V$) of all members, and the support reactions ($e_R$). The peak-response errors are $|\max\hat a-\max a|/\max a$ for the largest inter-storey drift ratio ($e_\mathrm{drift}$), the largest beam mid-span deflection relative to the chord ($e_\mathrm{defl}$; Hermite interpolation of the end rotations plus the fixed-end deflection $wL^4/(384EI)$), the largest absolute end moment ($e_{M,\max}$) and the largest absolute axial force ($e_{N,\max}$). The governing-member agreement is the fraction of frames in which the member with the largest end moment is predicted as the largest (top-1) or among the three largest (top-3). The equilibrium residual is $\|\mathbf{r}\|/\|\mathbf{P}_f-\mathbf{p}^0_f\|$ with $\mathbf{r}$ the nodal out-of-balance force at the free DOFs, moment rows divided by $H$; the compatibility residual is Eq.~\eqref{eq:compat}.

\paragraph{Statistics} The proposed ensemble is trained five times independently (five ensembles of five networks, 25 networks in total, with disjoint seeds). For each metric the per-run median and p95 are reported as mean $\pm$ standard deviation over the five runs, with a 95\% confidence interval from the $t$ distribution with four degrees of freedom. Within one run, 95\% confidence intervals of the median and p95 over test frames are obtained by the percentile bootstrap~\cite{efron1993bootstrap} with 2{,}000 resamples, and differences between two variants evaluated on the same frames (e.g.\ raw versus projected forces) by a paired bootstrap. Graph baselines are reported as mean $\pm$ standard deviation over three seeds and tabular baselines with bootstrap intervals.

\paragraph{Computing environment} All experiments ran on a cloud virtual machine with four cores of an Intel Xeon processor at 2.1\,GHz, 15\,GB of memory and no GPU, under Linux with Python 3.13, PyNite 3.2.0, NumPy 2.5, SciPy 1.18, scikit-learn 1.9, PyTorch 2.15 and OpenSeesPy 3.8. Networks were trained on the CPU with one thread each, four in parallel; timing measurements were taken separately on the otherwise idle machine (Section~\ref{sec:cost}).
